\clearpage
\thispagestyle{plain}
\section*{Vorwort}

\begin{center}
{\arabicfont\begin{Arabic}بِسْمِ اللَّهِ الرَّحْمَنِ الرَّحِيمِ\end{Arabic}}\\[3mm]
{\arabicfont\begin{Arabic}الْحَمْدُ لِلَّهِ، وَالصَّلَاةُ وَالسَّلَامُ عَلَى أَشْرَفِ الْأَنْبِيَاءِ وَالْمُرْسَلِينَ، نَبِيِّنَا مُحَمَّدٍ، وَعَلَى آلِهِ وَصَحْبِهِ أَجْمَعِينَ.\end{Arabic}}
\end{center}

\begin{quranvers}{Muḥammad ibn ʿAbd al-Wahhāb, al-Qawāʿid al-Arbaʿ, Einleitung}
{\arabicfont\begin{Arabic}أَسْأَلُ اللَّهَ الْكَرِيمَ رَبَّ الْعَرْشِ الْعَظِيمِ أَنْ يَتَوَلَّاكَ فِي الدُّنْيَا وَالْآخِرَةِ، وَأَنْ يَجْعَلَكَ مُبَارَكًا أَيْنَمَا كُنْتَ، وَأَنْ يَجْعَلَكَ مِمَّنْ إِذَا أُعْطِيَ شَكَرَ، وَإِذَا ابْتُلِيَ صَبَرَ، وَإِذَا أَذْنَبَ اسْتَغْفَرَ؛ فَإِنَّ هَؤُلَاءِ الثَّلَاثَ عُنْوَانُ السَّعَادَةِ.\end{Arabic}}
\end{quranvers}

Ich bitte Allah, den Großzügigen, den Herrn des gewaltigen Thrones, dass Er dein Schutzherr im
Diesseits und im Jenseits sei, dich segne, wo immer du auch bist, und dich zu jenen gehören
lasse, die bei einer Gabe dankbar sind, bei einer Prüfung geduldig bleiben und nach einer Sünde
um Vergebung bitten. Denn diese drei Dinge sind Zeichen der Glückseligkeit.

Wisse, mein lieber Bruder, dass zu den größten Wohltaten, die Allah \ar{ﷻ} einem Menschen
schenken kann, das Streben nach nützlichem Wissen gehört. Besonders in jungen Jahren ist es
eine große Gunst und Rechtleitung Allahs \ar{ﷻ}, wenn Er einem Menschen ermöglicht, seine Zeit
dafür zu nutzen, seine Religion zu erlernen und sich mit dem zu beschäftigen, was ihm im
Diesseits und im Jenseits nützt.

Gerade in unserer heutigen Zeit ist dies von besonderer Bedeutung. Viele Menschen sprechen
über die Religion, obwohl sie kein oder nur wenig Wissen besitzen. Jeder möchte seine Meinung
äußern und über Angelegenheiten sprechen, mit denen er sich möglicherweise nie ernsthaft
beschäftigt hat. Umso wichtiger ist es für dich, zu den Grundlagen zurückzukehren und deine
Religion auf Wissen und Beweisen aufzubauen.

Wissen ist jedoch kein Ziel, mit dem man sich schmückt oder sich über andere erhebt. Vielmehr
soll es dich Allah \ar{ﷻ} näherbringen, deine ʿAqīdah festigen und sich in deinen Worten und
Taten zeigen.

Dieses Buch soll dir dabei eine Grundlage geben. Es basiert auf dem Werk \emph{ad-Durūs
al-Muhimmah li-ʿĀmmati l-Ummah} von Shaykh ʿAbd al-ʿAzīz ibn Bāz \ar{رحمه الله}, in welchem
der Shaykh grundlegende Angelegenheiten zusammenfasste, die ein Muslim für seine Religion
benötigt.

Das Buch ist bewusst einfach gehalten und stellt keine umfassende Vertiefung in jede einzelne
Wissenschaft dar. Vielmehr soll es dir ein festes Fundament vermitteln, auf dem du – \ar{بإذن الله}
– weiter aufbauen kannst. Einzelne Angelegenheiten werden in den Erläuterungen ausführlicher
behandelt, doch auch diese stellen lediglich einen Einstieg in die jeweiligen Themen dar.

Betrachte dieses Buch daher nicht als Endstation, sondern als einen Anfang. Lerne, wiederhole
und festige das Gelernte. Frage nach, wenn du etwas nicht verstehst, und bemühe dich darum,
nach dem zu handeln, was du gelernt hast.

Und bitte Allah \ar{ﷻ} stets um nützliches Wissen, Aufrichtigkeit und Standhaftigkeit.

\vspace{1.5em}
\begin{flushright}
\textbf{Sahdjad Mahamadou}
\end{flushright}
\clearpage
